\section{Lista 2 (20/04/2026)}

\subsection{Exercícios obrigatórios}
\begin{exercise}
	\label{lista2:ex1}
	Construa variáveis aleatórias $X_1, \dots X_n$ tal que para todo
	$S \subsetneq [n]$,
	temos que $(X_i)_{i \in S}$ são independentes mas $X_1, \dots X_n$
	não são independentes.
\end{exercise}

Eu conheço um exemplo famoso relativamente simples. Seja
\[
	A = \{(v_1, \dots, v_n) \in \{0,1\}^n : \sum_{i = 1}^n v_i \equiv 0 \pmod 2\}
\]
e seja $(X_1, \dots, X_n)$ tomado uniformemente sobre $A$. Vamos ver
que $(X_i)$ satisfaz as condições do problema. Precisamos de um lema fácil.
\begin{lemma}
	Dado $n > 0 \in \Z$ e $c \in \{0,1\}$, seja
	\[ A_n = \{(v_1, \dots, v_n) \in \{0,1\}^n : \sum_{i = 1}^n v_i \equiv
	c \pmod 2\}\]
	então $|A_n| = 2^{n-1}$.
\end{lemma}

\begin{proof}
	Há formas diversas de fazer isso, por indução parece ser mais fácil.
	Claramente é verdade para $n = 1$.
	Suponha que vale para $n > 1$, para $n+1$ temos
	\begin{align*}
		|A_{n+1}| &= |A_{n+1}\cap\{v_{n+1} = 0\}| + |A_{n+1} \cap \{v_{n+1}=1\}|\\
		&= |\{(v_1,\dots,v_n) \in \{0,1\}^n : \sum_{i=0}^{n} v_i = c \pmod 2\}|\\
		&+ |\{(v_1,\dots,v_n) \in \{0,1\}^n : \sum_{i=0}^{n} v_i = c - 1 \pmod 2\}|\\
		&= 2^{n-1} + 2^{n-1} = 2^{n}
	\end{align*}
	como queríamos mostrar.
\end{proof}

\begin{proof}[ Prova do Exercício \ref{lista2:ex1}]
	Seja $A$ dado por
	\[
		A = \{(v_1, \dots, v_n) \in \{0,1\}^n : \sum_{i = 1}^n v_i \equiv 0 \pmod 2\},
	\]
	já sabemos que $|A| = 2^{n-1}$, tome $(X_1, \dots, X_n)$ um vetor
	uniformemente em $A$. Isso é, nossa variável aleatória $X_i$ é a
	i-ésima coordenada desse vetor aleatório.
	Vamos verificar que para todo $S \subsetneq [n]$, $(X_i)_{i \in S}$
	são independentes.
	Temos para toda sequência $(a_i)_{i \in S}$ de elementos $a_i \in \{0,1\}$ que
	\[
		\Pr(\forall i \in S,  X_i = a_i) = \frac{|\{\vec{v} \in \{0,
				1\}^{n-|S|}: \sum_{i=1}^{n-|S|} v_i \equiv c - \sum_{i \in S}
		a_i\}|}{2^{n-1}} = 2^{-|S|}
	\]
	pelo lema anterior. E, portanto, usando o lema novamente,
	\[
		\Pr(\forall i \in S,  X_i = a_i) = \prod_{i \in S}^{}\Pr(X_i = a_i) = 2^{-|S|}
	\]
	logo $(X_i)_{i \in S}$ são independentes (para ser preciso eu
		deveria considerar $\Pr(X_i \in B_i)$ onde $B_i \subset \{0,1\}$,
		mas é fácil ver que a equação anterior implica as igualdades dessa
	forma). No entanto, crucialmente,
	\[
		\Pr(X_n = 0 \; |\; X_1 = 0, X_2 = 0, \dots, X_{n-1} = 0) = 1 \neq
		\frac{1}{2} = \Pr(X_n = 0).
	\]
\end{proof}

\begin{exercise}
	Prove que
	\begin{enumerate}[label=(\alph*)]
		\item Para toda medida de probabilidade $\mu_X$ em $(\R,
			\mathcal{B}(\R))$ vale que $\mu_X(A) \in \{0,1\}$ para
			todo $A \in \mathcal{B}(\R)$ se e só se $\mu_X = \delta_p$ para
			algum $p \in \R$.
		\item Se $X$ é uma variável aleatória independente de si mesma,
			então $X = c$ quase certamente para $c \in \R$.
		\item Se $X$ é $c \in \R$ quase certamente, então $X$ é
			independente de qualquer outra variável aletória.
		\item Sejam $X$ e $Y$ independentes. Dê uma condição necessária e
			suficiente para $\Pr(X = Y) > 0$.
		\item Se $X$ e $Y$ são independentes tais que $X + Y = c$ quase
			certamente para algum $c \in \R$
			então $X = c_1$ e $Y = c_2$ quase certamente onde $c_1 + c_2 = c$.
		\item Prove que se $X$ e $Y$ são variáveis aleatórias independentes
			sem átomos, então $X + Y$ também não tem átomos.
	\end{enumerate}
\end{exercise}

Seguem as soluções.
\begin{enumerate}[label=(\alph*)]
	\item
		\begin{proof}
			A recíproca é óbvia, se $\mu = \delta_p$ então $\mu(A) = \id[p \in
			A]$. Para a ida, vamos encontrar
			um ponto com toda a massa. Note que se $A,B$ são disjuntos, não
			pode ser que $\mu(A) = \mu(B) = 1$.
			Portanto, para cada $n \in \N$, seja
			\[
				\R = \bigsqcup_{i \in \Z} \bigg[\frac{i}{2^n}, \frac{i+1}{2^n}\bigg)
			\]
			pela observação anterior, para cada $n$, existe somente um
			intervalo $A_n = [i/2_n, (i+1)/2^n)$
			com $\mu(A_n) = 1$. Claramente esses intervalos $A_n$ são
			encaixados e se intersectam em um ponto $p$.
			Vale que $\mu(\{p\}) = \lim_{n\to\infty} \mu(A_n) = 1$.
		\end{proof}

	\item
		\begin{proof}
			Para todo $A \in \mathcal{B}(\R)$, vale que
			\[
				\Pr(X \in A) = \Pr((X,X) \in A \times A) = \Pr(X \in A)^2,
			\]
			logo $\Pr(X \in A) = \{0,1\}$. Pela letra anterior,
			a medida associada a $X$ em $(\R, \mathcal{B}(\R))$ é uma $\delta$.
			Logo $X = c$ quase certamente
			para algum $c$.
		\end{proof}

	\item
		\begin{proof}
			Seja $X: \Omega \to \R$ quase certamente igual a $c \in \R$ e $Y:
			\Omega \to \R$ variável aleatória qualquer.
			Temos para $A \times B \in \mathcal{B}(\R)^2$ que
			\[
				\Pr((X,Y) \in A \times B) = \Pr(X^{-1}(A) \cap Y^{-1}(B)) = \id[c
				\in A] \cdot \Pr(Y^{-1}(B)) = \Pr(X \in A) \Pr(Y \in B)
			\]
			onde na segunda igualdade usamos que se $\mu(C) \in \{0,1\}$,
			então $\mu(C \cap D) = \mu(C)\mu(D)$. Segue
			que $(X,Y)$ são independentes na álgebra dos retângulos,
			aproximando indicadoras por retângulos ou usando Dynkin, segue o resultado.
		\end{proof}

	\item
		\begin{proof}
			Basta que exista $p \in \R$ com $\Pr(X = p) > 0$ e $\Pr(Y = p) >
			0$. Naturalmente,
			se tal $p$ existe então,
			\[
				\Pr(X = Y) \geq \Pr(X = Y = p) = \Pr(X=p) \Pr(Y = p) > 0.
			\]
			Se não existe tal $p$, então $X$ não possui átomos átomos em comum com $Y$.
			Sejam $(x_i)_{i \in B}$ os átomos de $X$ e $(y_i)_{i \in C}$ os
			átomos de $Y$ (sabemos que átomos sempre são enumeráveis).
			Logo, como $X,Y$ são independentes,
			\[
				\Pr(X = Y) = \int_{\R}\int_{\R} \id[x = y] p_X(x) p_Y(y) =
				\int_{\R} \Pr(X = y) p_Y(y)
			\]
			temos ainda
			\begin{align*}
				\int_{\R} \Pr(X = y) p_Y(y) &= \int_{s \not \in (y_i)_{i \in C}}
				\Pr(X = s) p_Y(s) + \sum_{i \in C} \Pr(X = y_i)\Pr(Y = y_i)\\
				&=\int_{s \not \in (y_i)_{i \in B}} \sum_{j \in B}\Pr(X =
				x_j)\id[s = x_j] p_Y(s) + \sum_{i \in C} \Pr(X = y_i)\Pr(Y = y_i)\\
				&=\sum_{j \in B} \Pr(X = x_j)\Pr(Y = x_j) + \sum_{i \in C} \Pr(X
				= y_i)\Pr(Y = y_i)\\
				&= 0
			\end{align*}

		\end{proof}

	\item
		\begin{proof}
			Se $X+Y = c$ q.c , então $X = c - Y$ q.c. Como $X$ é independente de $Y$,
			$X$ é independente de $X$, portanto, por (b), $X = c_1$ para algum
			$c_1 \in \R$ q.c.
			Por fim, temos $Y = c - c_1$ q.c.
		\end{proof}

	\item
		\begin{proof}
			Sendo $X,Y$ independentes, temos a distribuição produto e portanto,
			dado $c \in \R$,
			\[
				\Pr(X + Y = c) = \int_{\R}\int_\R \id[x + y = c] p_X(x)p_Y(y) =
				\int_{\R} \Pr(X = c - y) p_Y(y) = 0
			\]
			na verdade mostramos algo mais forte. Basta que $X$ não tenha átomos.
		\end{proof}
\end{enumerate}

\begin{exercise}
	Sejam $X_i \equiv \text{Exp}(1)$ i.i.d.
	\begin{enumerate}[label=(\alph*)]
		\item Prove que a.s
			\[ \limsup_n \frac{X_n}{\log n} = 1.\]
		\item Seja $M_n = \max_{i \in [n]} X_i$. Prove que a.s
			\[ \liminf_n \frac{M_n}{\log n} \geq 1.\]
		\item Prove que a.s. existe uma subsequência $n_k$ tal que
			\[ \limsup_k \frac{M_{n_k}}{\log n_k} \leq 1.\]
		\item Mostre que a.s.
			\[
				\lim_n \frac{M_n}{\log n} = 1.
			\]
	\end{enumerate}

\end{exercise}
Para esse exercício, enunciarei os lemas de Borel-Cantelli provados em aula.
\begin{lemma}
	\label{lemma:borel_cantelli_1}
	Sejam $(E_n)_{n \in \N}$ eventos em um espaço de probabilidade
	$(\Omega, \Sigma)$. Se
	\[
		\sum_{n \geq 1} \Pr(E_n) < \infty,
	\]
	então quase certamente somente finitos deles acontecem.
\end{lemma}

\begin{lemma}
	\label{lemma:borel_cantelli_2}
	Sejam $(E_n)_{n \in \N}$ eventos independentes espaço de
	probabilidade $(\Omega, \Sigma)$. Se
	\[
		\sum_{n \geq 1} \Pr(E_n) = \infty,
	\]
	então quase certamente infinitos deles acontecem.
\end{lemma}

\begin{enumerate}[label=(\alph*)]
	\item
		\begin{proof}
			Seja $\alpha > 1$ e $E_n = \{X_n/\log n > \alpha\}$. Vamos
			mostrar que a.s somente finitos $E_n$ ocorrem. Note que
			\[
				\sum_{n \geq 1} \Pr\bigg(\frac{X_n}{\log n} > \alpha\bigg) =
				\sum_{n \geq 1} \Pr(\text{Exp}(1) > \alpha \log n)
				= \sum_{n \geq 1} n^{-\alpha} < \infty
			\]
			pelo fato da última série ser convergente. Segue pelo lema
			[\ref{lemma:borel_cantelli_1}] que se $m \in \N$ e
			$A_m$ é o evento
			\[A_m := \bigg\{\frac{X_n}{\log n} > 1 + \frac{1}{m} \quad
				\text{finitas vezes}\bigg\}
			= \bigg\{\limsup_n \frac{X_n}{\log n} \leq 1 + \frac{1}{m}\bigg\}\]
			então $\Pr(A_m) = 1$. Temos por fim, $\Pr(\limsup_n X_n/\log(n)
			\leq 1) = \Pr(\bigcap_{m \in \N} A_m) = 1$.

			Falta a outra desigualdade, seja $\alpha < 1$ e defina os $E_n$
			igualmente. Note que os $E_n$ são independentes e
			crucialmente
			\[
				\sum_{n \geq 1} \Pr(E_n) = \sum_{n \geq 1} \Pr(\text{Exp}(1) >
				\alpha \log n) = \sum_{n \geq 1} n^{-\alpha} = \infty.
			\]
			Pelo lema [\ref{lemma:borel_cantelli_2}], para todo $\alpha < 1$,
			os $E_n$ acontecem infinitas vezes com probabilidade $1$.
			Ou seja, definindo
			\[
				B_m :=  \bigg\{\limsup_n \frac{X_n}{\log n} \geq 1 - \frac{1}{m}\bigg\}
			\]
			segue que $\Pr(B_m) = 1$ e portanto $\Pr(\limsup_n X_n/\log n
			\geq 1) = \Pr(\bigcap_{m \in \N} B_m) = 1$. Juntando as duas
			afirmações, segue que $\Pr(\limsup_n X_n/\log n = 1) = 1$.
		\end{proof}	

	\item
		\begin{proof}
			Seja $\alpha < 1$, vamos provar que $\Pr(\liminf_n M_n/\log n
			\geq \alpha) = 1$. Basta ver que a.s.
			somente para finitos termos vale que $M_n/\log n < \alpha$. Temos
			\[	
				\sum_{n \geq 1} \Pr(M_n /\log n < \alpha) = \sum_{n \geq 1}
				\Pr(\max_{i \in [n]} X_i \leq \alpha\log n) = \sum_{n \geq 1}
				\Pr( \bigcap_{i \in [n]} \{X_i \leq \alpha \log n\})
			\]
			segue por independência
			\[
				\sum_{n \geq 1} \Pr( \bigcap_{i \in [n]} \{X_i \leq \alpha \log
				n\}) = \sum_{n \geq 1} (1 - n^{-\alpha})^n \leq \sum_{n \geq 1}
				\exp(-n^{1 - \alpha}).
			\]
		\end{proof}
		Se $\alpha < 1$ essa série converge super rapidamente, já que
		$n^{1 - \alpha} \gg 2\log n$ e portanto existe $n_0$ tal que
		\[
			\sum_{n \geq n_0} \exp(-n^{1 - \alpha}) \leq \sum_{n \geq n_0}
			\exp(- 2\log n) = \sum_{n \geq n_0} \frac{1}{n^2} < \infty.
		\]
		Segue por Borel Cantelli [\ref{lemma:borel_cantelli_1}] que para
		todo tal $\alpha$, a.s. $Mn/ \log n < \alpha$ somente finitas vezes.
		Fazendo o mesmo argumento que antes (tomando interseção sobre uma
		sequência enumerável), temos $\Pr(\liminf M_n / \log n \geq 1) = 1$.

		Fiquei curioso em como fazer as próximas duas letras
		"corretamente", acredito que a minha solução não é a intencionada.
	\item[(c) e (d)]
		\begin{proof}
			Vamos ver que (a) implica $\limsup_n M_n/\log n \leq 1$ a.s. Por
			consequência (c) é verdade e juntando com (b), temos (d).
			Seja $\omega$ tal que $\limsup_n X_n(\omega)/\log n \leq 1$. Se
			$\alpha > 1$, existe $n_0(\omega) \in \N$ tal que para
			todo $n > n_0(\omega)$ vale
			\[
				X_n(\omega) \leq \alpha \log n.
			\]
			Agora, para esse $\omega$ fixo, se $n > n_0$, vale que
			\[
				\frac{M_n(\omega)}{\log n} = \max\bigg( \max_{j \leq n_0}
					\frac{X_j(\omega)}{\log n},  \max_{n_0 < m \leq
				n}\frac{X_m(\omega)}{\log n} \bigg)
			\]
			e $X_m/\log n \leq X_m/\log m \leq \alpha$ para $n_0 < m \leq n$.
			Portanto, tomando $\limsup$, temos
			\[
				\limsup_n \frac{M_n(\omega)}{\log n} \leq \max(\max_{j \leq n_0}
				\frac{X_j(\omega)}{\log n}, \alpha) = \alpha
			\]
			pois $X_j(\omega)/\log n \to 0$ quando $j$ está fixo. Portanto,
			para todo $\alpha > 1$,
			\[
				\{\limsup_n X_n/\log n \leq 1\} \subset \{\limsup_n M_n/\log n
				\leq \alpha\}
			\]
			ou seja, para todo tal $\alpha$,
			\[
				\Pr\{\limsup_n M_n/\log n \leq \alpha\} \geq \Pr\{\limsup_n
				X_n/\log n \leq 1\} = 1
			\]
			tomando interseções sobre enumeráveis, segue o resultado que
			\[
				\Pr\{\limsup_n M_n/\log n \leq 1\} = 1.
			\]
		\end{proof}
\end{enumerate}

\begin{exercise}
	Seja $(X_n)_n$ uma sequência de variáveis aleatórias e
	defina $\mathcal{B}_n = \sigma((X_k)_{k \geq n})$ e
	$\mathcal{B}_\infty = \bigcap_n \mathcal{B}_n$.
	\begin{enumerate}[label=(\alph*)]
		\item Considere $f : \N \to \R_+$. Mostre que os seguintes eventos
			estão em $\mathcal{B}_\infty$.
			\begin{enumerate}[label=(\roman*)]
				\item $\{X_n/f(n) \in B \ \text{i.o}\}$ para $B \in \mathcal{B}(\R)$.
				\item $\{\limsup_n X_n/f(n) > c\}$ para $c \in \R$.
				\item $\{\exists \lim_n X_n/f(n)\}$.
			\end{enumerate}
		\item Conclua que se $\mathcal{E} := \{\exists \lim_n X_n/f(n)\}$,
			então $\limsup_n X_n/f(n)$ e $\lim_n \id[\mathcal{E}]X_n/f(n)$ são
			variáveis aleatórias $\mathcal{B}_\infty$ mensuráveis.
	\end{enumerate}
\end{exercise}
Seguem as soluções.
\begin{enumerate}[label=(\alph*)]
	\item
		\begin{proof}
			Vamos escrever cada um dos conjuntos como uma interseção esperta.
			Para (i), temos
			\[
				\{X_n/f(n) \in B \ \text{i.o}\} = \bigcap_{n \geq 1} \bigcup_{m
				\geq n} \{X_m/f(m) \in B\}
			\]
			como para todo $m \geq n$, $\{X_m/f(m) \in B\} \in \mathcal{B}_n$,
			segue para todo $k \in \N$,
			\[
				\bigcap_{n \geq 1} \bigcup_{m \geq n} \{X_m/f(m) \in B\} =
				\bigcap_{n \geq k} \bigcup_{m \geq n} \{X_m/f(m) \in B\} \in \mathcal{B}_k
			\]
			ou seja $\{X_n/f(n) \in B \ \text{i.o}\} \in \bigcap_k
			\mathcal{B}_k = \mathcal{B}_\infty$.

			Para (ii) fazemos algo muito similar,
			abrindo a definição
			\[
				\{\limsup_n X_n/f(n) > c\} \quad = \quad \{\exists\ d > c,
				\ X_n/f(n) \in (d,\infty) \ \text{i.o}\}.
			\]
			Por (i), para cada $d \in \R$ o evento $\{X_n/f(n) \in (d,\infty)
			\ \text{i.o}\} \in \mathcal{B}_\infty$. Portanto,
			\[
				\bigcup_{d \in \mathbb{Q}, d > c} \{X_n/f(n) \in (d,\infty)
				\ \text{i.o}\} \in \mathcal{B}_\infty
			\]
			e isso é equivalente ao que queríamos mostrar.

			Em (iii) podemos expressar a existência do limite como a sequência
			ser de Cauchy. Notamos que
			\[
				\{\exists \lim_n X_n/f(n)\} = \{X_n/f(n) \ \text{é de Cauchy}\} =
				\bigcap_{\substack{\eps \in \mathbb{Q},\\ \eps > 0}} \bigcup_{n
				\geq 1} \bigcap_{m \geq n} \{|X_n/f(n) - X_m/f(m)| < \eps\}.
			\]			
			Segue assim como antes (brincando com conjuntos crescentes e
			decrescentes) que para todo $k \in \N$,
			\[
				\bigcup_{n > 1} \bigcap_{m \geq n} \{|X_n/f(n) - X_m/f(m)| <
				\eps\} = \bigcup_{n > k} \bigcap_{m \geq n} \{|X_n/f(n) -
				X_m/f(m)| < \eps\} \in \mathcal{B}_k
			\]
			e portanto,
			\[
				\bigcup_{n > 1}\bigcap_{m \geq n} \{|X_n/f(n) - X_m/f(m)| <
				\eps\} \in \bigcap_{k \in \N} \mathcal{B}_k = \mathcal{B}_\infty.
			\]
			Fazendo a interseção para $\eps \in \mathbb{Q}$ com $\eps > 0$,
			segue o resultado $\{\exists \lim_n X_n/f(n)\} \in \mathcal{B}_\infty$.
		\end{proof}

	\item
		\begin{proof}
			Segue diretamente da letra (a) que $\id(\mathcal{E})$ é mensurável
			e, notando que os conjuntos $(c,\infty)$ com $c \in \R$ geram
			$\mathcal{B}(\R)$, que
			$\limsup_n X_n/f(n)$ é mensurável. Mas então fica claro que
			\[
				\lim_n \id[\mathcal{E}]X_n/f(n) = \id[\mathcal{E}] \limsup_n
				X_n/f(n) \in \mathcal{B}_\infty.
			\]
		\end{proof}
\end{enumerate}

% \subsection{Outros exercícios}

% \begin{exercise}
% 	(2) Seja $X = (X_1, \dots, X_d)$ um vetor gaussiano, i.e. $X$ tem densidade
% 	\[
% 		p_X(x) = (2\pi)^{d/2}|\det(K)|^{-1/2}\exp\bigg(\frac{- (x - \mu)^T
% 		K^{-1} (x - \mu)}{2}\bigg)
% 	\]
% 	onde $\mu \in \R^d$ e $K$ é uma matriz positiva definida. Mostre que
% 	$(X_1, \dots X_d)$ são independentes
% 	se e somente se $\text{Cov}(X_i, X_j) = 0$ para todo $i \neq j$.
% \end{exercise}
% \begin{proof}
% 	Na verdade, vale que $K$ é a matriz de covariância das coordenadas.
% 	(Esse é meio chato de fazer a conta).
% \end{proof}

\pagebreak
